% ##############################################################################
\section{Evaluation}
\label{sec:evaluation}

We evaluate \uiPolicy{} along five dimensions: truthfulness, operational leverage, policy depth, runtime cost, and defense in depth. The results are from one production deployment of \system{}. They should be read as an experience report, not a controlled experiment.

\subsection{Threat model}
\label{sec:threat_model}

\uiPolicy{} addresses stale UI gates, impersonation abuse, and client-side authorization bypass. It does not attempt to replace standard cloud controls for secrets, network isolation, IAM, or database security.

The policy store is the new critical surface. Moving policy from code to DynamoDB shifts the blast radius of a bad write from a code branch to the UI contract of a tenant. The mitigation is not denial. Operators must treat policy-store writes with the same seriousness as code deploys.

\subsection{Truthful UX coverage}
\label{sec:truthful_coverage}

Table~\ref{tab:eval_summary} summarizes the production policy surface and coverage. Coverage means that a surface conditionally renders based on at least one CASL ability check, module membership test, or module-layout field.

\begin{table}[H]
  \centering
  \small
  \begin{tabularx}{\textwidth}{X r}
    \toprule
    \textbf{Metric} & \textbf{Value} \\
    \midrule
    Modules defined in policy store & 7 \\
    CASL permission subjects & 25+ \\
    \code{moduleLayout} config fields & 150+ \\
    Tenant onboardings requiring deployment & 0 \\
    \midrule
    \textbf{Policy surface coverage} & \\
    Navigation items & 8/11 (73\%) \\
    Module components & 7/7 (100\%) \\
    Configuration tabs & 5/6 (83\%) \\
    \midrule
    \textbf{UI elements hidden for restricted tenants} & \\
    Single-module tenant & 6/7 (86\%) \\
    Two-module tenant & 5/7 (71\%) \\
    \bottomrule
  \end{tabularx}
  \caption{Policy adoption and coverage in production. The remaining navigation entries are generic surfaces such as profile, documentation, and global account actions.}
  \label{tab:eval_summary}
\end{table}

High-value surfaces are covered. Module components are fully policy-driven. Configuration tabs are almost fully covered. Navigation still has universal entries, which is expected. Those entries are intentionally shared across tenants and roles.

Before adoption, permission and navigation mismatches produced recurring support escalations that needed developer investigation. Since adoption, we have observed no drift incidents against policy-gated surfaces. This is not a universal claim about all authorization bugs. It is a narrower result: the architecture eliminated the stale-gate class for surfaces that were migrated to the contract.

\subsection{Runtime cost}
\label{sec:runtime_cost}

Table~\ref{tab:runtime} reports full-materialization latency for \code{/v2/me}. The path includes JWT validation, tenant and role reads, layout-fragment assembly, impersonation-list assembly, and CASL rule construction.

\begin{table}[H]
  \centering
  \small
  \begin{tabularx}{0.55\textwidth}{X r}
    \toprule
    \textbf{Metric} & \textbf{512 MB Lambda (ms)} \\
    \midrule
    P50 & 856 \\
    P75 & 938 \\
    P90 & 1,038 \\
    P95 & 1,112 \\
    P99 & 1,176 \\
    \bottomrule
  \end{tabularx}
  \caption{Controller Lambda full-materialization latency from 1,368 production observations, March--April 2026, 512 MB configuration.}
  \label{tab:runtime}
\end{table}

Full materialization is not the steady-state path. ETag-conditional revalidations dominate normal traffic at a 91\% cache-hit rate. Those return HTTP 304 with a P50 of 68 ms, with median user-visible revalidation latency of roughly 75 ms.

\subsection{Operational leverage}
\label{sec:operational_leverage}

The production deployment serves more than ten tenants, each with multiple users across roles. Adding a tenant requires a policy-store record and Cognito group assignment. Enabling a module or changing layout is a policy update. No code deployment is required for routine tenant onboarding.

The operational gain is mundane and valuable. A tenant change stops being a front-end release. Support can reason from one contract. QA can fetch \code{/v2/me} and inspect exactly which product world the user will see.

\subsection{Defense in depth}
\label{sec:defense_in_depth}

Six layers protect the path between browser and backend.

\begin{table}[H]
  \centering
  \small
  \begin{tabularx}{\textwidth}{l c X}
    \toprule
    \textbf{Layer} & \textbf{Where enforced} & \textbf{Mechanism} \\
    \midrule
    Edge middleware & Server & Session check, path validation, and security headers. \\
    Server components & Server & Redirect on missing or invalid session. \\
    CASL abilities & Server + client & Built server-side from \code{/v2/me}; exported for UX only. \\
    Module gating & Server & \code{can('view', module)} before layout is returned. \\
    View-as signing & Server & HMAC-signed effective context; secret never reaches the browser. \\
    Control plane & Server & JWT validation, allow-list membership, and schema-prefix derivation. \\
    \bottomrule
  \end{tabularx}
  \caption{Defense-in-depth layers. The client participates in UX decisions, but backend and control-plane checks remain authoritative.}
  \label{tab:defense_layers}
\end{table}

A broken client-side rule should at worst produce a confusing interface. It should not produce a permitted backend response.

\subsection{What the evaluation does not show}
\label{sec:eval_limits}

The evaluation is correctness-oriented. It shows that covered surfaces derive from a server contract and that the deployed system carries useful production metrics. It does not include a controlled before/after incident study, nor does it prove absence of drift in legacy surfaces. Longitudinal drift measurement and a full tenant-matrix CI gate remain future work.
